% Generated draft: requires mathematical and source review.
\documentclass[11pt]{article}
\usepackage[margin=25mm]{geometry}
\usepackage{fontspec}
\setmainfont{Latin Modern Roman}
\usepackage{amsmath,amssymb,mathtools,unicode-math}
\setmathfont{Latin Modern Math}
\usepackage{enumitem,needspace,xurl,hyperref}
\setlength{\parindent}{0pt}
\setlength{\parskip}{5pt}
\setlength{\emergencystretch}{3em}
\newcommand{\N}{\mathbb{N}}
\newcommand{\Z}{\mathbb{Z}}
\newcommand{\Q}{\mathbb{Q}}
\newcommand{\R}{\mathbb{R}}
\newcommand{\C}{\mathbb{C}}
\newcommand{\F}{\mathbb{F}}
\newcommand{\A}{\mathbb{A}}
\newcommand{\PP}{\mathbb{P}}
\newcommand{\E}{\mathbb{E}}
\newcommand{\eps}{\varepsilon}
\newcommand{\dd}{\,\mathrm{d}}
\begin{document}
\section*{UnsolvedMath ID: 800951}
\textbf{Source problem code:} catalogue entry
\subsection{Definitions and mathematical statement}
Let \(M\) be a closed surface of genus at least \(2\), as in the supplied source, and let \(S^1=\R/\Z\). Let \(g\) be a Riemannian metric on \(M\times S^1\) with no conjugate points: along every geodesic, a Jacobi field that vanishes at two distinct parameter values is identically zero. Define
\[
p:M\times\R\longrightarrow M\times S^1,\qquad p(x,t)=(x,t+\Z),
\qquad \widetilde g=p^*g.
\]
This is the infinite cyclic cover obtained by unwrapping the circle factor; its covering subgroup is \(\pi_1(M)\times\{0\}\subseteq\pi_1(M)\times\Z\).

The supplied Conjecture 31, attributed in context to Burago--Kleiner, asserts that for every such \(M\) and \(g\), the lifted metric splits as a Riemannian product. Interpreted as an isometric splitting, the assertion is that there exist a Riemannian metric \(g_1\) on \(M\) and a diffeomorphism \(F:M\times\R\to M\times\R\) such that
\[
F^*(g_1\oplus\dd t^2)=\widetilde g.
\]
The source writes an equality of products without specifying whether the splitting must respect the original product coordinates or any prescribed deck action. The formulation above expresses an isometric splitting; these additional compatibility requirements remain unresolved. The source also does not explicitly specify orientability in its surface hypothesis.
\subsection{Short English statement}
The conjecture asks whether every Riemannian metric without conjugate points on the product of a closed surface of genus at least \(2\) and a circle becomes isometric to a surface times a Euclidean line after unwrapping the circle. The supplied statement leaves the required compatibility of this splitting with the original coordinates unspecified.
\subsection{Sources}
\begin{itemize}
\item Supplied local UnsolvedMath registry record: permanent ID \(800951\), exact source code \texttt{catalogue entry}, in the AMR Open Problem Lists collection. The record labels the problem open; this is a supplied status, not an independently verified mathematical or literature assessment.
\item Burns and Matveev, \emph{Open Problems and Questions about Geodesics}, source-list year supplied as 2019, Conjecture 31. Supplied source URL: \url{https://arxiv.org/abs/1308.5417}. The registry lists Burns and Matveev as proposers and gives Burago--Kleiner as attribution/context.
\item Supplied extraction provenance: \texttt{arxiv\_latex}; access date recorded as 2026-10-16. This date and provenance are supplied metadata; no independent source visit or verification is claimed. Supplied upstream revision: \texttt{d371ee917ac9647fbb06cd6fc529deec6f3dda82}.
\end{itemize}
\end{document}
